\documentclass[11pt,aspectratio=169]{beamer}

\usepackage[T2A]{fontenc}
\usepackage[utf8]{inputenc}
\usepackage[english,russian]{babel}
\usepackage{lmodern}
\usepackage{amsmath,amssymb}
\usepackage{booktabs}
\usepackage{xcolor}
\usepackage{tikz}
\usetikzlibrary{arrows.meta,positioning,shapes.geometric}

\usetheme{Boadilla}
\usecolortheme{seahorse}

\definecolor{primary}{HTML}{1F3A5F}
\definecolor{accent}{HTML}{C9302C}
\definecolor{soft}{HTML}{4A6FA5}
\definecolor{paper}{HTML}{F8F8F4}
\definecolor{good}{HTML}{1E6F3E}
\definecolor{bad}{HTML}{B23A48}

\setbeamercolor{structure}{fg=primary}
\setbeamercolor{frametitle}{bg=primary,fg=white}
\setbeamercolor{title}{fg=primary}
\setbeamercolor{block title}{bg=primary,fg=white}
\setbeamercolor{block body}{bg=paper}
\setbeamertemplate{navigation symbols}{}

\title[Leveraged carry на Pendle]{%
  Доходная стратегия на фиксированной\\ доходности Pendle с защитой от риска}
\subtitle{Финальный проект курса <<Блокчейн и DeFi>>}
\author{Михаил Швайков}
\date{Май 2026}

\begin{document}

% ============================================================
\begin{frame}[plain]
  \titlepage
\end{frame}

\begin{frame}{План доклада}
  \begin{enumerate}\itemsep0.5em
    \item Задача и идея стратегии: что мы строим и почему
    \item Алгоритм действий и гипотеза работоспособности
    \item Математика стратегии: доходность, газ, комиссии
    \item Защита от ликвидации: адаптивный контроллер
    \item Хеджирование: активное (Hyperliquid) и пассивное (Pendle Boros)
    \item Эмпирическая проверка: 6~циклов на 3~сетях
    \item Программная реализация и интеграция с fractal-defi
    \item Тестирование, ограничения, выводы
  \end{enumerate}
\end{frame}

% ============================================================
\section{Идея}

\begin{frame}{Что мы строим}
  Pendle разделяет любой токен с переменной доходностью (например sUSDe от Ethena)
  на две части:
  \begin{itemize}\itemsep0.2em
    \item \textbf{PT} (Principal Token) --- обещание выплатить 1 единицу
          базового токена в день погашения. До этой даты торгуется со
          скидкой, и величина скидки и есть \emph{фиксированная} доходность.
    \item \textbf{YT} (Yield Token) --- забирает плавающую часть доходности.
  \end{itemize}

  \vspace{0.4em}
  Мы работаем только с PT. Идея: купить PT (зафиксировав хорошую доходность),
  заложить в Morpho, занять стейблкоин под этот залог, на занятое снова купить
  PT. Повторив 5 раз, получаем плечо $\approx$3{,}36$\times$ на разнице между
  доходностью PT и ставкой займа.

  \vspace{0.4em}
  Дополнительно к этому держим параллельную позицию, зарабатывающую на
  ставке финансирования (funding rate) бессрочных фьючерсов --- это
  второй независимый источник доходности и одновременно частичная защита
  от риска роста implied yield.
\end{frame}

\begin{frame}{Алгоритм действий: пошагово}
  \footnotesize
  \begin{enumerate}\itemsep0.3em
    \item Покупаем PT-sUSDE на Pendle AMM за стартовые \$10\,000~USDC.
    \item Кладём купленные PT в Morpho как залог.
    \item Берём заём в USDC, размер выбираем так, чтобы отношение долга к
          залогу составило $L = 0{,}80$ (запас 11{,}5~п.п. до порога автоматической
          ликвидации $L_{\text{liq}} = 0{,}915$).
    \item На занятые USDC покупаем ещё PT.
    \item Повторяем шаги 2--4 ещё четыре раза.
    \item Параллельно открываем хеджирующую позицию (long на ставку
          финансирования ETH-перпа), размер $= 1\times$ от итогового
          залога в PT.
    \item Держим всю позицию до даты погашения PT. Раз в час проверяем
          фактическое плечо: если оно вышло за безопасный коридор --- ребалансируем.
    \item В день погашения PT автоматически меняется на базовый токен 1:1.
          Гасим долг, закрываем хедж, фиксируем результат.
  \end{enumerate}
\end{frame}

\begin{frame}{Гипотеза работоспособности}
  \begin{block}{Тезис}
    Премия leveraged-стратегии над пассивным удержанием PT существует.
    Она компенсирует риск принудительной ликвидации. Грамотное управление
    плечом сохраняет премию, а параллельный хедж добавляет второй
    источник доходности и одновременно частично компенсирует риск.
  \end{block}

  \vspace{0.4em}
  \textbf{Что нужно проверить эмпирически:}
  \begin{itemize}\itemsep0.2em
    \item Действительно ли стратегия прибыльна после учёта реальных
          комиссий и газа?
    \item Достаточен ли запас по плечу 11{,}5~п.п. в наблюдаемых
          рыночных режимах?
    \item Работает ли хедж как обещано, и какова его реальная
          доходность?
    \item Что происходит в стрессовый период?
  \end{itemize}

  \vspace{0.4em}
  Дальше расскажу, как мы это проверяли.
\end{frame}

% ============================================================
\section{Математика}

\begin{frame}{Базовая идея в одной картинке}
  \footnotesize
  \centering
  \begin{tikzpicture}[
      node distance=2.2cm and 1.6cm,
      box/.style={draw=primary, very thick, rounded corners=4pt,
                  minimum height=1.0cm, minimum width=2.6cm,
                  align=center, fill=paper},
      arr/.style={-{Stealth[length=3mm]}, very thick, primary}
    ]
    \node[box] (s) {Стартовые\\ \$10\,000};
    \node[box, right=of s] (pt) {Купили PT\\ на \$10\,000};
    \node[box, right=of pt] (m) {Заложили в\\ Morpho};
    \node[box, below=of m] (b) {Заняли 80\%\\ от стоимости};
    \node[box, below=of pt] (p2) {Купили PT\\ на занятое};
    \node[box, below=of s] (c) {Повторяем\\ 5 циклов};
    \draw[arr] (s) -- (pt);
    \draw[arr] (pt) -- (m);
    \draw[arr] (m) -- (b);
    \draw[arr] (b) -- (p2);
    \draw[arr] (p2) -- (c);
    \draw[arr] (c) -- (s);
  \end{tikzpicture}

  \vspace{0.5em}
  После 5 циклов на \$10\,000 капитала у нас PT-токенов на \$33\,600 ---
  плечо \textbf{3{,}36$\times$}. Долг при этом \$26\,900.
\end{frame}

\begin{frame}{Откуда доходность: чистая формула}
  Пусть $L$ --- доля долга от залога. После 5 циклов:
  \[
    \text{Залог в PT}: \;C_5 = X \cdot \frac{1-L^5}{1-L}, \qquad
    \text{Долг}: \;D_5 = LX \cdot \frac{1-L^5}{1-L}.
  \]

  Чистая годовая доходность на стартовый капитал (без учёта затрат):
  \[
    \boxed{\;\text{APY} \;\approx\;
      \frac{r_{PT} - L \cdot r_b}{1-L} \cdot \frac{1-L^5}{5(1-L)}\;}
  \]

  где $r_{PT}$ --- доходность PT, $r_b$ --- ставка займа.

  \vspace{0.4em}
  \textbf{Пример (реальный цикл sUSDE-27NOV2025, Ethereum):}\\
  средняя $r_{PT} = 6{,}68\%$, $r_b = 5{,}62\%$, $L = 0{,}80$:\\
  APY $\approx (0{,}0668 - 0{,}8 \cdot 0{,}0562)/0{,}20 \cdot 3{,}36/5 \approx 7{,}3\%$ годовых.

  \vspace{0.3em}
  По остальным циклам спред $r_{PT} - L \cdot r_b$ колеблется
  от \textbf{2{,}0\%} (sUSDE-5FEB2026) до \textbf{10{,}5\%}
  (стрессовый weETH-27JUN2024) --- отсюда разброс конечной доходности.

  \vspace{0.4em}
  \textbf{Это теоретический потолок.} Дальше учтём всё, что съедает доходность.
\end{frame}

\begin{frame}{Учёт комиссии Pendle AMM}
  \footnotesize
  Каждая покупка/продажа PT идёт через автоматический маркет-мейкер Pendle.
  Стоимость складывается из двух частей:

  \vspace{0.3em}
  \begin{block}{Комиссия пула (fee)}
    Pendle берёт \textbf{0{,}1\%} от входной суммы каждого свопа.
    На 5 циклов открытия плюс 5 на закрытие это \textbf{1\%}
    от исходного капитала.
  \end{block}

  \vspace{0.2em}
  \begin{block}{Проскальзывание (slippage)}
    Цена сдвигается из-за размера сделки относительно глубины пула:
    $\text{slip} = \sigma \cdot \text{размер сделки}/\text{ликвидность пула}$,
    где $\sigma = 0{,}5$ --- консервативный коэффициент (сделка
    размером с весь пул двигает цену на 50\%).
  \end{block}

  \vspace{0.3em}
  \normalsize
  Эти затраты вычитаются внутри модели PT-сущности на каждом свопе,
  не как фиксированная величина в конце цикла.
\end{frame}

\begin{frame}{Учёт газа: реалистичные цифры по сетям}
  Один цикл стратегии = 3 транзакции (своп Pendle $\approx$400k газа +
  залог Morpho $\approx$150k + заём Morpho $\approx$250k). При 10 gwei и
  ETH = \$2\,500 это даёт оценку:

  \vspace{0.4em}
  \footnotesize
  \begin{tabular}{lrrr}
    \toprule
    Сеть & Открытие позиции & Закрытие & Одна ребалансировка \\
    \midrule
    Ethereum (главная)  & \$150 & \$125 & \$30 \\
    Arbitrum            &  \$10 &   \$8 &  \$2 \\
    Base                &   \$8 &   \$6 & \$1{,}5 \\
    \bottomrule
  \end{tabular}

  \vspace{0.7em}
  \normalsize
  На Ethereum полный круг = \$275 за газ. На капитале \$10\,000 это сразу
  $-2{,}75\%$ от стартового баланса.

  \vspace{0.4em}
  \textbf{Это центральное наблюдение работы.} Газ съедает теоретическую
  премию плеча на главной сети при мелком капитале. Поэтому либо нужен
  больший капитал (от \$50\,000+), либо переход на L2, либо дополнительный
  доходный поток --- хедж.

  \vspace{0.4em}
  В коде газ моделируется как фиксированное списание из стартового
  капитала перед началом цикла. Это математически эквивалентно оплате
  газа из того же кошелька.
\end{frame}

% ============================================================
\section{Защита от ликвидации}

\begin{frame}{Что такое ликвидация и почему она опасна}
  Если цена PT упадёт настолько, что отношение долга к залогу
  превысит порог $L_{\text{liq}} = 0{,}915$ --- Morpho \emph{мгновенно}
  продаёт залог по любой доступной цене, чтобы погасить долг. Разница
  ложится на нас как убыток.

  \vspace{0.4em}
  \textbf{Просадка $\ne$ ликвидация.} Просадка --- временное падение
  оценки, отыгрывается к погашению PT (мы получаем номинал). Ликвидация ---
  принудительная фиксация убытка прямо сейчас, при этом мы теряем не только
  падение цены, но и комиссию ликвидатора (5--7{,}5\% от залога).

  \vspace{0.4em}
  Поэтому стратегия должна следить за плечом и реагировать \emph{до}
  того, как оно достигнет порога. Это и есть задача адаптивного контроллера.
\end{frame}

\begin{frame}{Адаптивный контроллер: как работает}
  \footnotesize
  Каждый час алгоритм проверяет фактическое плечо $L_t = D_t/C_t$.

  \vspace{0.4em}
  \normalsize
  \textbf{Верхняя граница безопасного коридора} $L_U$ ищется так:
  \emph{какое максимальное плечо я могу держать, чтобы вероятность
  ликвидации в ближайшие 3 часа не превысила $10^{-4}$ (один шанс
  на десять тысяч)?}

  \vspace{0.4em}
  Из теории броуновского движения с дрейфом (формула Бахелье):
  \[
    \Pr(\text{ликвидация за } h) \;=\;
    \Phi\!\left(\frac{b - \mu h}{\sigma \sqrt{h}}\right)
    + e^{2\mu b / \sigma^2}\,
      \Phi\!\left(\frac{b + \mu h}{\sigma \sqrt{h}}\right)
  \]
  где $b = \ln(L/L_{\text{liq}})$ --- расстояние до порога в логарифмах,
  $\mu, \sigma$ --- наблюдаемые дрейф и волатильность цены PT,
  $\Phi$ --- стандартная нормальная функция распределения.

  \vspace{0.4em}
  Решаем это уравнение численно (бисекция) и получаем $L_U$.
  \textbf{Если $L_t > L_U$ --- продаём часть PT, гасим часть долга, опускаем плечо до целевого.}
\end{frame}

\begin{frame}{Кто исполняет ребалансировку в реальности}
  \footnotesize
  \vspace{0.3em}
  \begin{block}{В нашей работе}
    Алгоритм реализован как off-chain бэктестер на Python, который
    обрабатывает исторические данные. Это доказательство
    работоспособности модели, не готовый продукт.
  \end{block}

  \vspace{0.3em}
  \begin{block}{Как бы это работало в продакшене}
    \begin{itemize}\itemsep0.15em
      \item \textbf{Smart-contract vault} (Yearn-style): пользователь
            кладёт деньги, контракт держит позицию, off-chain бот
            (например, Gelato, Chainlink Automation) подаёт
            ребаланс-транзакции.
      \item \textbf{DIY}: пользователь сам пишет бота, который раз в
            час смотрит state и подаёт транзакцию со своего адреса.
    \end{itemize}
  \end{block}

  \vspace{0.3em}
  \normalsize
  В обоих случаях транзакцию подаёт \emph{внешний агент}, не
  смарт-контракт сам. Morpho-контракт лишь \emph{выполняет}
  пришедшую инструкцию.
\end{frame}

% ============================================================
\section{Хеджирование}

\begin{frame}{Зачем хедж: экономическая цепочка}
  Доходность sUSDe формируется из ставки финансирования (funding rate),
  которую Ethena получает, держа короткие ETH-перпы на биржах.

  \vspace{0.4em}
  \textbf{Что происходит, когда funding rate растёт:}
  \begin{enumerate}\itemsep0.2em
    \item Ethena зарабатывает больше $\Rightarrow$ sUSDe доходность растёт
    \item Рынок ждёт большую доходность от PT-sUSDE $\Rightarrow$
          implied yield этого PT растёт
    \item Цена PT-токена падает (mark-to-market), наше плечо растёт
  \end{enumerate}

  \vspace{0.4em}
  \begin{block}{Идея хеджа}
    Купить позицию, которая зарабатывает на росте funding rate.
    Тогда: рынок плохой для PT $\to$ хорош для хеджа. И наоборот:
    при positive funding (нормальный режим) хедж приносит carry
    параллельно с PT.
  \end{block}

  \vspace{0.4em}
  Хедж работает в обе стороны: и как страховка, и как второй
  источник доходности.
\end{frame}

\begin{frame}{Pendle Boros: что это}
  Boros --- отдельный продукт Pendle Finance, запущенный в августе
  2025~года. Он токенизирует \emph{форвардный контракт} на среднюю
  ставку финансирования конкретного перпетуала (например, ETH-USDC
  на Hyperliquid) до фиксированной даты.

  \vspace{0.4em}
  \textbf{Как работает.} Каждый Boros-рынок --- это собственный AMM,
  торгующий <<подразумеваемой средней ставкой>> (mark APR). Длинная
  позиция в Boros получает фактический funding минус подразумеваемый,
  выплаченный при покупке. По экономике это аналогично активному
  шортингу перпа, но \emph{без управления маржой}: купил токен,
  держишь до экспирации, получаешь усреднённый поток.

  \vspace{0.4em}
  \begin{block}{Почему именно Boros для пассивной версии стратегии}
    Один токен в кошельке вместо позиции на бирже деривативов с
    margin management. Подходит для retail-инвестора, который не
    хочет ежечасно следить за здоровьем своей позиции на Hyperliquid.
  \end{block}
\end{frame}

\begin{frame}{Два варианта хеджа в нашей работе}
  \footnotesize
  \begin{tabular}{p{0.46\linewidth}p{0.46\linewidth}}
    \toprule
    \textbf{Активный (Hyperliquid)} & \textbf{Пассивный (Pendle Boros)} \\
    \midrule
    Сами держим short ETH-перпа на Hyperliquid, получаем funding payments. &
    Покупаем Boros-токен, держим до экспирации. \\
    \addlinespace[2pt]
    Высокая операционная сложность: маржа, ликвидация, ребалансировка. &
    Низкая: купил и забыл до даты погашения. \\
    \addlinespace[2pt]
    Получаем актуальную funding rate (волатильную). &
    Получаем форвардное ожидание ставки (сглаженное). \\
    \addlinespace[2pt]
    Исторические данные доступны за 5+ лет. &
    Публичный API даёт только последние 3 недели (запущен Aug 2025). \\
    \bottomrule
  \end{tabular}

  \vspace{0.5em}
  \normalsize
  Поскольку для Boros нет полных исторических данных, мы построили
  \textbf{синтетический эквивалент}: умножили Hyperliquid funding на
  коэффициент $0{,}45$ (эмпирически измеренное соотношение Boros/HL
  на свежем окне в 21 день).

  \vspace{0.3em}
  Это даёт реалистичную \emph{нижнюю границу} доходности для пассивного
  инвестора через Boros. Активная версия через Hyperliquid даёт
  верхнюю границу.
\end{frame}

% ============================================================
\section{Эмпирическая проверка}

\begin{frame}{6 циклов, 3 сети}
  \footnotesize
  \begin{tabular}{p{3.4cm}p{1.5cm}p{0.7cm}p{7.5cm}}
    \toprule
    Цикл & Сеть & Дней & Характеристика \\
    \midrule
    sUSDE-25SEP2025 & Ethereum & 60 & лето-осень 2025, спокойный \\
    sUSDE-27NOV2025 & Ethereum & 59 & осень 2025, спокойный \\
    sUSDE-5FEB2026  & Ethereum & 60 & зима 2025/26, сжатие доходностей \\
    USDe-11DEC2025  & Arbitrum & 60 & скачок implied yield, уход ликвидности \\
    USDe-11DEC2025  & Base     & 60 & то же, ещё тоньше пул \\
    \textbf{weETH-27JUN2024} & \textbf{Ethereum} & \textbf{60} &
    \textbf{стрессовый: LRT-кризис апреля 2024} \\
    \bottomrule
  \end{tabular}

  \vspace{0.5em}
  \normalsize
  Стартовый капитал во всех экспериментах \$10\,000. Шаг наблюдений 1 час.
  Стрессовый цикл выбран как период максимальной волатильности PT за
  всю доступную историю.

  \vspace{0.3em}
  \footnotesize
  \textbf{LRT} (liquid restaking token) --- токенизированный stETH с
  дополнительной доходностью от перестейкинга на EigenLayer.
  eETH от ether.fi --- крупнейший LRT. В апреле 2024 весь сегмент
  испытал кризис: eETH временно потерял прижатость к ETH, implied yield
  PT-weETH скакнул.
\end{frame}

\begin{frame}{Откуда берутся данные}
  \footnotesize
  \begin{tabular}{p{3.6cm}p{10.8cm}}
    \toprule
    \textbf{Источник} & \textbf{Что и зачем} \\
    \midrule
    Pendle REST API & Цена PT, implied yield, ликвидность пула, время
    до погашения. Основной канал для PT-метрики. \\
    \addlinespace[2pt]
    Morpho GraphQL API & Ставка займа, утилизация пула, параметр LLTV.
    Используется при моделировании долга. \\
    \addlinespace[2pt]
    Binance public REST & Цена sUSDe/USDe относительно USDT.
    Используется при погашении PT (если базовый токен депегировал,
    получим меньше \$1 за каждый PT). \\
    \addlinespace[2pt]
    Hyperliquid info API & Историческая funding rate ETH-перпетуала ---
    модель активного хеджа. \\
    \addlinespace[2pt]
    Boros REST API & Текущий mark APR Boros-токенов --- только для
    валидации соотношения Boros/Hyperliquid. Исторических данных
    мало для полноценного бэктеста. \\
    \bottomrule
  \end{tabular}

  \vspace{0.4em}
  \normalsize
  Все пять источников --- публичные, не требуют API-ключей. Это
  важно: работа полностью воспроизводима без процедуры регистрации
  и оплаты подписок.
\end{frame}

\begin{frame}{Главные результаты на спокойных циклах}
  \footnotesize
  \begin{tabular}{lrrrr}
    \toprule
    Цикл & PT без плеча & Static & Hedged-Active & Hedged-Boros \\
    \midrule
    sUSDE-25SEP2025 & $+12{,}2\%$ & \textcolor{bad}{$+2{,}8\%$}  & \textcolor{good}{$+38{,}6\%$} & $+18{,}9\%$ \\
    sUSDE-27NOV2025 &  $+7{,}1\%$ & \textcolor{bad}{$-0{,}7\%$} & \textcolor{good}{$+24{,}2\%$} & $+10{,}5\%$ \\
    sUSDE-5FEB2026  &  $+5{,}1\%$ & \textcolor{bad}{$-3{,}5\%$} & \textcolor{good}{$+19{,}8\%$} &  $+7{,}0\%$ \\
    \bottomrule
  \end{tabular}

  \vspace{0.7em}
  \normalsize
  Все цифры --- годовая доходность на стартовый капитал \$10\,000 на главной
  сети Ethereum, с учётом газа \$275 и комиссий Pendle.

  \vspace{0.4em}
  \textbf{Два вывода:}
  \begin{itemize}\itemsep0.2em
    \item Простой плечевой loop на главной сети при стартовом капитале
          \$10\,000 \textbf{неприбылен}: газ съедает премию.
    \item Хеджирование \emph{всегда} переводит стратегию в плюс ---
          и пассивное (Boros) даёт $+7$\,до\,$+19\%$ годовых, и активное
          (Hyperliquid) даёт $+20$\,до\,$+39\%$.
  \end{itemize}
\end{frame}

\begin{frame}{Что произошло на стрессовом цикле}
  \footnotesize
  \textbf{Контекст.} В апреле 2024 eETH временно потерял прижатость
  к ETH (LRT-кризис: накопилось много рискованных позиций с плечом
  на restaking). PT-weETH, который привязан к eETH, торговался по
  цене $\sim$0{,}949 при моменте нашего <<входа>> --- это implied yield
  около 30\% годовых.

  \vspace{0.6em}
  \normalsize
  \begin{tabular}{lrr}
    \toprule
    Стратегия & Итог за 60 дней & Год.\,доходность \\
    \midrule
    PT без плеча       & \$10\,522 & $+31{,}8\%$ \\
    Static loop        & \$11\,021 & $+62{,}2\%$ \\
    Hedged-Active      & \$12\,130 & $+129{,}8\%$ \\
    Hedged-Boros       & \$11\,520 & $+92{,}6\%$ \\
    \bottomrule
  \end{tabular}

  \vspace{0.6em}
  \textbf{Важное уточнение.} Это \emph{не} <<хедж спас от просадки>>.
  PT-токены сами заработали много: мы купили дёшево (implied yield
  на нашем входе уже был высоким), к погашению цена вернулась к 1.
  Параллельно ETH-funding в апреле--июне был сильно положительным,
  и хедж независимо заработал свой carry.

  \vspace{0.3em}
  То есть это \textbf{двойная независимая прибыль}, а не страховка.
\end{frame}

\begin{frame}{Аномалия циклов USDe на L2}
  \footnotesize
  На Arbitrum и Base тот же цикл (USDe-11DEC2025) показал убытки
  у всех стратегий. Это \emph{не} ошибка кода, а реальное событие.
  Причины:

  \vspace{0.5em}
  \normalsize
  \begin{enumerate}\itemsep0.2em
    \item \textbf{Резкое падение ликвидности пула}: TVL рухнул со \$163k
          до \$29k. Большие свопы стали дорогими (проскальзывание 5--10\%
          от размера сделки).
    \item \textbf{Скачок implied yield} в последние дни цикла: $6{,}6\% \to 30{,}6\%$
          --- но \emph{не благоприятный} для нас, потому что мы вошли
          \emph{при низком} implied yield. Цены PT падали для нас
          mark-to-market, и к погашению восстановились слабо.
    \item \textbf{На L2 нет рынков для sUSDe} --- только сам USDe.
          Базовая доходность ниже.
  \end{enumerate}

  \vspace{0.5em}
  \footnotesize
  \textbf{Важно различать со стресс-циклом.} Там мы \emph{вошли}
  при высоком implied yield (купили дёшево). Здесь мы \emph{вошли}
  при низком (купили дорого), и рост implied yield ударил по нам.
  Знак входа всё решает.

  \vspace{0.4em}
  \normalsize
  \textbf{Вывод.} Главная сеть Ethereum для PT-sUSDE остаётся
  оптимальной площадкой, несмотря на дорогой газ: пул толще,
  underlying-доходность выше.
\end{frame}

\begin{frame}{Адаптивное плечо: ни разу не активировалось}
  Во всех шести проверенных циклах фактическое плечо ни разу не
  вышло за верхнюю границу безопасного коридора. Кривая капитала
  адаптивной стратегии \emph{идентична} простой до цента.

  \vspace{0.5em}
  \begin{block}{Это не провал, а подтверждение конструкции}
    \large
    Адаптивное управление \emph{бесплатно} в спокойных и умеренно
    стрессовых режимах. Контроллер активируется только на грани
    реальной ликвидации.
  \end{block}

  \vspace{0.5em}
  Проверочный эксперимент с агрессивным плечом $L = 0{,}84$
  (буфер 7{,}5~п.п. вместо 11{,}5): контроллер активируется немедленно,
  стоит $\sim$60 базисных пунктов как страховая премия.

  \vspace{0.4em}
  \textbf{Эмпирическая проверка защитного механизма не проведена.}
  В наших циклах не было события, которое бы заставило контроллер
  сработать в реалистичной конфигурации. Это известное ограничение
  работы.
\end{frame}

% ============================================================
\section{Программная реализация}

\begin{frame}{Что построено: общий обзор}
  Работа реализована как полноценный Python-пакет на базе библиотеки
  \texttt{fractal-defi} (учебный фреймворк курса). Покрытие 153 тестами,
  публикация под лицензией BSD-3.

  \vspace{0.4em}
  \begin{block}{Три уровня}
    \footnotesize
    \begin{itemize}\itemsep0.2em
      \item \textbf{Сущности} (entities): типизированные модели протоколов
            с собственным состоянием. Что-то меняется --- например, при
            <<покупке PT>> --- состояние сущности корректно обновляется.
      \item \textbf{Загрузчики} (loaders): обёртки над публичными API
            пяти источников данных. Кэшируют ответы на диск, чтобы
            бэктесты выполнялись быстро.
      \item \textbf{Стратегии}: логика, которая каждый час смотрит на
            текущее состояние сущностей и решает, какие действия
            запросить (например, <<занять 7440 USDC>>).
    \end{itemize}
  \end{block}
\end{frame}

\begin{frame}{Что написано для fractal-defi}
  \texttt{fractal-defi} в исходной версии (1.3.1) имел сущности для
  Aave, Hyperliquid, Uniswap V2/V3, Lido stETH и GMX. \textbf{Pendle и
  Morpho полностью отсутствовали.} Мы добавили:

  \vspace{0.4em}
  \footnotesize
  \begin{tabular}{p{3.6cm}p{10.8cm}}
    \toprule
    Класс & Что делает \\
    \midrule
    \texttt{PendlePTEntity} & Принципал-токен Pendle: своп через AMM
    с комиссией и проскальзыванием, обработка экспирации, погашение
    PT с учётом возможного депега базового токена. \\
    \addlinespace[2pt]
    \texttt{MorphoEntity} & Изолированный рынок Morpho Blue: начисление
    процентов на долг, контроль LLTV, флаг ликвидации. \\
    \addlinespace[2pt]
    \texttt{FundingHedgeEntity} & Длинная позиция на ставку финансирования
    с линейным начислением PnL. \\
    \bottomrule
  \end{tabular}

  \vspace{0.6em}
  \normalsize
  Эти три класса -- вклад в библиотеку.
  Может использоваться любым другим исследователем для своих стратегий
  на Pendle или Morpho.
\end{frame}

\begin{frame}{Скрипты для воспроизведения работы}
  \footnotesize
  \begin{tabular}{p{4.4cm}p{10.0cm}}
    \toprule
    Скрипт & Что делает \\
    \midrule
    \texttt{run\_backtest.py} & Прогон бэктеста на одном цикле, выводит
    таблицу метрик. \\
    \addlinespace[2pt]
    \texttt{multi\_cycle\_backtest.py} & Прогон шести циклов на трёх
    сетях, статистическая валидация. \\
    \addlinespace[2pt]
    \texttt{sweep\_hedge\_ratio.py} & Сетка по размеру хеджа от 0 до 1{,}5
    --- чувствительность к выбору размера хеджа. \\
    \addlinespace[2pt]
    \texttt{validate\_boros\_proxy.py} & Сравнение Boros mark APR с
    Hyperliquid funding на свежих 21 днях. \\
    \addlinespace[2pt]
    \texttt{build\_observations.py} & Подготовка наблюдений из всех
    источников в одном проходе. \\
    \bottomrule
  \end{tabular}

  \vspace{0.5em}
  \normalsize
  Любой из этих скриптов запускается одной командой, без настройки
  API-ключей и переменных окружения.
\end{frame}

% ============================================================
\section{Тестирование}

\begin{frame}{Как мы убедились, что стратегия действительно работает}
  Один эмпирический бэктест на одном цикле --- слабое доказательство.
  Мы применили четыре техники из современной методологии validation
  алгоритмических стратегий.

  \vspace{0.4em}
  \begin{enumerate}\itemsep0.3em
    \item \textbf{Множественные циклы.} Шесть независимых временных окон
          в трёх сетях покрывают разные рыночные режимы:
          спокойный (3 цикла), периоды роста implied yield (2 цикла),
          стрессовый LRT-кризис (1 цикл).
    \item \textbf{Стресс-тест волатильности.} Искусственно увеличиваем
          наблюдаемую волатильность PT в 1{,}5 раза и пересчитываем ---
          проверяем, что адаптивный контроллер реагирует адекватно
          даже на не виданное в данных движение.
    \item \textbf{Block-bootstrap.} Ресэмплируем почасовые доходности
          блоками по 24 часа, реконструируем кривую капитала. Это даёт
          доверительные интервалы.
    \item \textbf{Out-of-sample.} Калибруем параметры контроллера на
          одном цикле, применяем на другом --- проверяем, что они
          переносимы.
  \end{enumerate}
\end{frame}

\begin{frame}{Block-bootstrap: разница статистически значима}
  Цикл sUSDE-27NOV2025. 200 бутстрэп-траекторий, блоки по 24 часа:

  \vspace{0.5em}
  \begin{center}
  \begin{tabular}{lrrr}
    \toprule
    Стратегия & $p_5$ & медиана & $p_{95}$ \\
    \midrule
    Static loop          & \$10\,047 & \$10\,186 & \$10\,339 \\
    Hedged-Active (HL)   & \$10\,435 & \$10\,605 & \$10\,797 \\
    \bottomrule
  \end{tabular}
  \end{center}

  \vspace{0.6em}
  Нижний 5\% защищённой стратегии (\$10\,435) выше верхнего 95\%
  простой (\$10\,339). Доверительные интервалы не пересекаются ---
  превосходство защищённой стратегии \textbf{не случайное наблюдение},
  а статистически значимый эффект.

  \vspace{0.5em}
  Проверка формулы первого касания методом Монте-Карло
  (30\,000 траекторий $\times$ 2\,400 шагов): расхождение менее $1\%$.
  Аналитика согласуется с симуляцией.
\end{frame}

\begin{frame}{Почему стратегия не сломается на реальных данных}
  \footnotesize
  \begin{enumerate}\itemsep0.3em
    \item \textbf{Газ и комиссии моделируются по реальным сетевым
          параметрам.} Это не теоретический бэктест --- если стратегия
          прошла у нас, на реальной сети получим похожий результат.
    \item \textbf{Шесть циклов на трёх сетях} --- невозможно подобрать
          параметры одновременно под все. Если бы переподогнали под
          один цикл, остальные показали бы другие числа.
    \item \textbf{Стрессовый цикл eETH-кризиса} --- максимальная
          волатильность PT за всю доступную историю. Если стратегия
          прошла его, обычные рыночные колебания её не сломают.
    \item \textbf{Известны граничные условия}: на главной сети при
          мелком капитале неприбыльна, на L2 при тонких пулах
          неприбыльна. Это не сюрприз, а \emph{понятые ограничения
          применимости}.
    \item \textbf{Адаптивный контроллер} построен на математической
          модели первого касания, проверенной симуляцией. В нормальном
          режиме PT-цена ведёт себя как броуновский процесс ---
          предположения модели согласуются с данными.
  \end{enumerate}
\end{frame}

% ============================================================
\section{Итоги}

\begin{frame}{Ограничения работы (важные)}
  \begin{enumerate}\itemsep0.3em
    \item \textbf{Эмпирически не проверили срабатывание адаптивного
          контроллера в реальной угрозе ликвидации.} В нашей выборке
          такого события не случилось. Контроллер построен на правильной
          математической модели, но эмпирического подтверждения
          защитного механизма в работе нет.
    \item \textbf{Boros и Hyperliquid funding --- разные инструменты.}
          Корреляция 0{,}47, средний разрыв 584 базисных пункта на
          свежем окне. Поэтому даём \emph{диапазон} ожидаемой
          доходности: верхняя граница через Hyperliquid, нижняя через
          синтетический Boros, истинное значение между ними.
    \item \textbf{Не учли риск смарт-контрактов.} В стеке три протокола
          (Pendle, Morpho, Ethena); риск потерь от
          эксплойтов.
    \item \textbf{Размер выборки.} Шесть циклов --- небольшая статистика;
          мы смягчили её block-bootstrap-ом.
  \end{enumerate}
\end{frame}

\begin{frame}{Главные выводы}
  \footnotesize
  \begin{enumerate}\itemsep0.3em
    \item \textbf{Стратегия работает, но не везде и не всегда.}
          Простая версия на главной сети с мелким капиталом
          неприбыльна после газа. Это важная честная находка.

    \item \textbf{Хедж экономически необходим} на главной сети ---
          не косметика, а условие прибыльности. Активный (Hyperliquid):
          $+20$\,до\,$+39\%$ годовых, пассивный (Boros): $+7$\,до\,$+19\%$.

    \item \textbf{В стрессовом режиме} стратегия не <<защитилась от
          просадки>>, а заработала на двух независимых линиях ---
          фиксированной доходности PT при удачном входе плюс carry
          положительного funding rate.

    \item \textbf{Адаптивный контроллер} построен правильно, но в наших
          данных не понадобился. Эмпирическая проверка защитного
          механизма --- задача будущей работы.

    \item \textbf{Реальное применение} осмысленно для капитала
          от \$50\,000 с активным управлением, или пассивно через
          Boros. Это не стратегия для retail с мелкими суммами.
  \end{enumerate}
\end{frame}

\begin{frame}[plain]
  \centering
  \vspace{1.5cm}
  {\LARGE \textbf{Спасибо за внимание}}\\[2em]
  {\Large Готов ответить на вопросы}

  \vspace{2cm}
  {\footnotesize
  \color{soft}
  Код:~~\texttt{github.com/qqmikhasik/pendle-pt-loop}\\
  Whitepaper:~~8 страниц, в репозитории}
\end{frame}

\end{document}
